\documentclass[11pt,a4paper]{article}
\usepackage[utf8]{inputenc} % Not needed with XeLaTeX
\usepackage{fontspec}
\usepackage[margin=0.75in]{geometry}
\usepackage{xcolor}
\definecolor{darknavy}{RGB}{0,0,100}
\usepackage[colorlinks=true, urlcolor=darknavy, linkcolor=darknavy]{hyperref}
\usepackage{enumitem}
% \usepackage{fontawesome}
\usepackage{fontawesome5}
\usepackage{titlesec}
\usepackage{academicons}

% % --- Font: Source Sans Pro via Type1 (works with pdfLaTeX)
% \usepackage[default,scale=1.0]{sourcesanspro} % body + sans by default
% \renewcommand{\familydefault}{\sfdefault}      % make sans the default family
% \usepackage[scaled]{helvet}
\usepackage[default,scale=0.95]{opensans}

% (optional) math that pairs reasonably with sans body
\usepackage{newtxsf} % sans math; or use newtxtext/newtxmath for serif math

\titleformat{\section}{\Large\bfseries}{\thesection}{1em}{}[\titlerule]
\titlespacing{\section}{0pt}{12pt}{6pt}

\usepackage{comicneue}

\begin{document}

\begin{center}
  % {\Huge\textbf{Pavel Popov}}\\[0.5em]
  % {\comicneue\Huge\textbf{Pavel Popov}}\\[0.5em]
  {\fontspec{JetBrains Mono}\Huge Pavel Popov}\\[0.5em]
  {\fontspec{JetBrains Mono}\footnotesize +1 (470) 920-1876 \quad | \quad Center street NW, Atlanta, GA 30318} \quad | \quad {paavali.popov@gmail.com}\\
  {\small \href{https://github.com/paavalipopov}{Github: paavalipopov} \quad | \quad
  % \href{https://www.linkedin.com/in/pavel-popov-83b69a123/}{\faLinkedin\ Pavel Popov} \quad
  \href{https://www.linkedin.com/in/pavel-popov-83b69a123/}{LinkedIn: pavel-popov-83b69a123} \quad | \quad
  \href{https://scholar.google.com/citations?user=4alUD5UAAAAJ}{Google scholar: 4alUD5UAAAAJ}}
\end{center}

\section*{Skills}
\begin{itemize}
[leftmargin=*,noitemsep]
\item \textbf{Machine Learning:} Deep learning architecture design, transformer architectures, representation learning, time series analysis, explainable/interpretable AI.
\item \textbf{Large-Scale Training:} Distributed data parallel \& multi-GPU training, HPC \& SLURM clusters, model optimization at scale.
\item \textbf{Media \& Signal Modeling:} Large-scale imaging \& time-series (video) data, digital signal processing, feature engineering.
\item \textbf{Programming \& Tools:} Python, PyTorch, Scikit-learn, Pandas, MATLAB, Git, Docker, SQL, mongoDB, Spark.
\end{itemize}
\vspace{-18pt}
\section*{Open Source Projects}
\noindent\textbf{\href{https://pypi.org/project/ml4fmri/}{ml4fmri}: Python toolkit for training \& benchmarking DL architectures on large-scale time-series data}\\
[-2em]
\begin{itemize}[leftmargin=*,noitemsep]
    \item Unified 13 DL architectures (including transformers) for multivariate time series analysis under a standardized training/evaluation API.
    \item Automated an end-to-end training \& benchmarking pipeline (ingestion, training, evaluation) for reproducible large-scale experiments.
\end{itemize}

\noindent\textbf{\href{https://pypi.org/project/brainbow/}{brainbow} --- Visualization tool for multi-dimensional neuroimaging data}\\
[-2em]
\begin{itemize}[leftmargin=*,noitemsep]
    \item Built a command-line visualization tool for rendering and exploring 3D/4D neuroimaging data formats.
\end{itemize}
\vspace{-18pt}
\section*{Experience}
\noindent{\textbf{Center for Translational Research in Neuroimaging and Data Science}} \hfill January 2022 --- Now\\
Graduate Research Assistant \hfill \textit{Atlanta, USA}\\
[-2em]
\begin{itemize}[leftmargin=*,noitemsep]
    \item Designed and trained deep neural architectures (including transformers) for large-scale multimodal imaging \& time-series data (4D fMRI --- volumetric image/video data).
    \item Developed a robust DL classifier for MRI image analysis with SOTA performance and 10x speedup over existing benchmarks.
    \item Scaled model training across multi-GPU HPC clusters with distributed data parallel training (SLURM), optimizing throughput on large imaging datasets.
    \item Investigated representation learning and model interpretability to understand semantic representations learned by deep models on imaging \& time-series data.
    \item Collaborated with cross-functional research teams to translate findings into shared tooling; presented findings at internal meetings and international conferences.
\end{itemize}

\noindent\textbf{Institute of Radioengineering and Electronics} \hfill May 2017 --- September 2020\\
Research Assistant \hfill \textit{Moscow, Russia}\\
[-2em]
\begin{itemize}[leftmargin=*,noitemsep]
    \item Modeled complex signal propagation and high-frequency signal dynamics in novel hardware architectures based on wave logic.
    \item Developed mathematical signal-processing models providing theoretical validation for prototype hardware controllers.
    \item Ran large-scale numerical simulations to validate theoretical models.
\end{itemize}
\vspace{-18pt}
\section*{Education}
\noindent\textbf{Georgia State University} \hfill June 2027\\
PhD in Computer Science \hfill
% Machine Learning \& Neuroimaging,
GPA 3.89\\[0.3em]
\textbf{Moscow Institute of Physics and Technology} \hfill June 2020\\
MS in Applied Mathematics and Physics \hfill
% Spintronics,
GPA 3.92\\[0.3em]
\textbf{Moscow Institute of Physics and Technology} \hfill June 2018\\
BS in Applied Mathematics and Physics \hfill
% Solid State Physics,
GPA 3.81

\end{document}
